% Auto-generated by scripts/evaluate_report_metrics.py
% Recommended packages: \usepackage{booktabs} \usepackage{amsmath} \usepackage{amssymb}

\subsection{技术指标与评估方法}
\label{sec:metrics-eval}

\paragraph{指标口径总览}
表\,\ref{tab:metrics} 汇总了平台的关键技术指标。除“组件库规模/模块数”等可直接统计的指标外，其余指标均通过可复现实验流程测得，并在下文给出数据集、判定规则与统计口径。

\paragraph{数据集与环境}
我们构造了一个覆盖多场景的需求集合用于评估（中英文混合，包含少量歧义与缺省表述以避免指标虚高）。需求解析评估样本数为 $N=41$；组件选择评估样本数为 $N=25$。为保证可复现性，本次评估采用离线模式运行（不调用外部 LLM API）。实验在本地开发机上完成（OS: Windows-10-10.0.26200-SP0; Python: 3.10.10）。

\paragraph{说明（避免指标被误读）}
本节“需求解析准确率”评估的是\textbf{离线启发式回退解析器}在\textbf{约束明确陈述}的样本上对关键字段的抽取正确性，用于验证系统在无外部 LLM 依赖时仍可交付与可复现；该指标\textbf{不等同于}在线 LLM 模式下对任意自由表述的泛化能力。在线模式可在同一评估集上复跑以获得对比结果。

\subsubsection{需求解析准确率}
需求解析输出包含多字段（方案类型、目标平台、安全等级、是否要求后量子等）。我们采用“字段级准确率”的平均作为总体解析准确率：
\begin{equation}
\mathrm{Acc}_\mathrm{parse} = \frac{1}{|F|} \sum_{f\in F} \frac{1}{N}\sum_{i=1}^N \mathbb{I}\left[\hat{y}_i^{(f)} = y_i^{(f)}\right],
\label{eq:acc-parse}
\end{equation}
其中 $F$ 表示被评估字段集合，$\mathbb{I}[\cdot]$ 为指示函数。该口径可以避免“某一个字段错误就整条判错”导致的过严统计，同时便于定位误差来源。

\begin{table}[htbp]
  \centering
  \caption{需求解析字段级正确性（$N=41$，离线回退；Wilson 95\% CI）}
  \label{tab:parse-acc}
  \begin{tabular}{lrrrr}
    \toprule
    字段 & 正确/总数 & 通过率(\%) & 95\% CI 下限 & 95\% CI 上限 \\
    \midrule
    方案类型 & 41/41 & 100.00 & 91.43 & 100.00 \\
    目标平台 & 41/41 & 100.00 & 91.43 & 100.00 \\
    安全等级 & 41/41 & 100.00 & 91.43 & 100.00 \\
    后量子需求 & 41/41 & 100.00 & 91.43 & 100.00 \\
    \midrule
    字段平均（$\mathrm{Acc}_\mathrm{parse}$） & -- & 100.00 & -- & -- \\
    严格全字段同时正确 & 41/41 & 100.00 & 91.43 & 100.00 \\
    \bottomrule
  \end{tabular}
\end{table}

\noindent\textit{注：}当出现“41/41”这类点估计为 $100\%$ 的结果时，我们同时报告 Wilson 95\% 置信区间以反映样本量不确定性；例如严格全字段通过率的 95\% CI 下限为 91.43\%，可作为更保守的表述依据。

\subsubsection{组件选择正确率}
对每条需求，我们生成 $K$ 个候选方案（本文默认 $K=3$）。当候选集中至少存在一个方案同时满足：1) 方案类型与需求一致；2) 覆盖需求中明确声明的偏好组件（如 AES+GCM、ChaCha20-Poly1305、Kyber-768 等）；3) 若需求声明后量子安全，则候选中需包含至少一种 PQC 组件（如 Kyber/Dilithium/ML-\*），则记为一次“命中”。组件选择正确率定义为命中比例：
\begin{equation}
\mathrm{Acc}_\mathrm{select} = \frac{1}{N}\sum_{i=1}^N \mathbb{I}\left[\exists\, s \in \mathcal{S}_i: s\ \text{is acceptable}\right].
\label{eq:acc-select}
\end{equation}

本次评估命中 22/25，即 88.00\%，Wilson 95\% CI 为 [70.04, 95.83]。为避免指标虚高，我们额外加入了少量“超出当前启发式覆盖范围”的偏好样本（例如 Camellia-GCM、RSA-4096、SM4-GCM），用于体现系统在“偏好不在库/不在候选集”时的边界表现；本次未命中样本为：S23, S24, S25。

\subsubsection{方案生成时延（P95）}
我们统计从提交需求到输出最终交付包（候选方案、审计结果与导出产物索引）的端到端耗时，并报告 $95\%$ 分位数（P95）：
\begin{equation}
T_{95} = \mathrm{Percentile}(\{t_i\}_{i=1}^N, 95\%).
\label{eq:t95}
\end{equation}

\begin{table}[htbp]
  \centering
  \caption{端到端时延统计（离线可复现模式，$N=41$）}
  \label{tab:latency}
  \begin{tabular}{lc}
    \toprule
    统计量 & 时延(ms) \\
    \midrule
    平均值 & 36.9 \\
    P95 & 39.8 \\
    最大值 & 46.7 \\
    \bottomrule
  \end{tabular}
\end{table}

\subsubsection{性能仿真误差}
我们选择典型算法组合（AES-GCM、ChaCha20-Poly1305、SHA-256），在本机用标准密码库执行真实加密/哈希操作测得吞吐，作为参考值；再与平台性能估计器输出对比。误差采用平均绝对百分比误差（MAPE）：
\begin{equation}
\mathrm{Err}_\mathrm{perf} = \frac{1}{M}\sum_{j=1}^M \left| \frac{\hat{p}_j - p_j}{p_j} \right| \times 100\%,
\label{eq:err-perf}
\end{equation}
其中 $p_j$ 为参考吞吐，$\hat{p}_j$ 为估计吞吐。

\begin{table}[htbp]
  \centering
  \caption{吞吐估计对比（平台: desktop，payload=1MB）}
  \label{tab:perf}
  \begin{tabular}{lrrr}
    \toprule
    算法 & 实测吞吐(MB/s) & 估计吞吐(MB/s) & 绝对百分比误差(\%) \\
    \midrule
    AES-GCM-128 & 2836.90 & 2200.00 & 22.45\\
    AES-GCM-256 & 2635.37 & 1600.00 & 39.29\\
    ChaCha20-Poly1305 & 1679.62 & 1725.00 & 2.70\\
    SHA-256 & 2304.11 & 1380.00 & 40.11\\
    \midrule
    平均绝对百分比误差（MAPE） & \multicolumn{3}{r}{26.14\%} \\
    \bottomrule
  \end{tabular}
\end{table}

\subsubsection{漏洞模式识别率与 CVE 识别率}
漏洞检测分为两类：1) 弱模式/不安全配置（如 MD5、SHA-1、ECB、静态 IV 等）的规则识别；2) 已知 CVE 的特征匹配识别。分别统计在测试集合上命中预期类别的比例，作为“漏洞模式识别率”与“CVE 识别率”。

\begin{table}[htbp]
  \centering
  \caption{漏洞检测评估结果（Wilson 95\% CI）}
  \label{tab:vuln}
  \begin{tabular}{lrrrr}
    \toprule
    任务 & 命中/样本 & 命中率(\%) & 95\% CI 下限 & 95\% CI 上限 \\
    \midrule
    弱模式/不安全配置识别 & 10/10 & 100 & 72.25 & 100.00 \\
    CVE 特征匹配识别 & 17/17 & 100 & 81.57 & 100.00 \\
    \bottomrule
  \end{tabular}
\end{table}

\begin{table}[htbp]
  \centering
  \caption{技术指标汇总（本次评估结果）}
  \label{tab:metrics}
  \begin{tabular}{lll}
    \toprule
    指标类别 & 指标名称 & 实测值 \\
    \midrule
    功能完整性 & 组件库规模 & 152 个组件（primitives 122 / modes 16 / protocols 14） \\
    功能完整性 & 功能模块数 & 19 个模块（4 个角色 + 15 个工具服务） \\
    准确性 & 需求解析准确率（离线回退） & 100.00\%\ (41/41; 95\%CI [91.43\%, 100.00\%]) \\
    准确性 & 组件选择正确率 & 88.00\%\ (22/25; 95\%CI [70.04\%, 95.83\%]) \\
    性能 & 方案生成时延（P95） & $<40\,ms$ \\
    性能 & 性能仿真误差（MAPE） & $\pm$ 26.14\% \\
    安全性 & 漏洞模式识别率 & 100\%\ (10/10; 95\%CI [72.25\%, 100.00\%]) \\
    安全性 & CVE 识别率 & 100\%\ (17/17; 95\%CI [81.57\%, 100.00\%]) \\
    \bottomrule
  \end{tabular}
\end{table}
